%% Currículo de exemplo - [YOUR_NAME]
%% Visão geral de competências, experiência e realizações.
%% Use como referência-mestra ao personalizar currículos por vaga.
%%
%% SETUP: rode /setup para preencher com seus dados reais.
%% Compile com: lualatex main_example.tex
%% (pdflatex costuma falhar em MiKTeX moderno com erros de expansão de fonte do fontawesome5.)

\documentclass[11pt,a4paper,sans]{moderncv}
\moderncvstyle{banking}
\moderncvcolor{blue}

% Força o nome e os títulos de seção a renderizarem no azul do moderncv (color1).
% O estilo banking padrão os deixa pretos: o \colorlet do moderncvstylebanking.sty
% copia (não referencia) a cor de destaque anterior ao esquema, então as cores do
% nome são congeladas antes do \moderncvcolor rodar. Reatribua depois. \namefont é
% o hook por onde toda macro de estilo de nome passa, então isto também funciona no
% moderncv 2.3.1 (apt do Debian/Ubuntu), que não tem \firstnamestyle/\lastnamestyle.
\renewcommand*{\namefont}{\fontsize{34}{36}\bfseries\upshape}
\colorlet{firstnamecolor}{color1}
\colorlet{lastnamecolor}{color1}
\colorlet{namecolor}{color1}
\renewcommand*{\sectionstyle}[1]{{\sectionfont\color{color1}#1}}

\usepackage[utf8]{inputenc}
% Hifenização em português brasileiro. Sem isso o LaTeX hifeniza português com
% padrões do inglês e quebra palavras em lugares errados - mais visível no resumo
% profissional justificado, e custa linhas num currículo com limite duro de 2
% páginas. Remova (ou troque a opção) se o currículo não for em português.
\usepackage[brazil]{babel}
% O moderncv carrega o hyperref por conta própria num hook \AtEndPreamble, então o
% \hypersetup precisa ir num \AtEndPreamble nosso: no moderncv < 2.4 um
% \usepackage{hyperref} de topo conflita com o \RequirePackage[unicode]{hyperref}
% da própria classe. A partir da 2.4.0 a classe passa as opções via
% \PassOptionsToPackage, que é o que elimina esse conflito.
\AtEndPreamble{\hypersetup{
    colorlinks=true,
    linkcolor=blue,
    filecolor=magenta,
    urlcolor=blue,
    pdftitle={[YOUR_NAME] - Currículo},
    % Mantenha pdfpagemode=UseNone: este bloco roda depois do \AtEndPreamble do
    % próprio moderncv (moderncv.cls define pdfpagemode lá), então um valor
    % FullScreen aqui venceria e abriria todo currículo em modo apresentação.
    pdfpagemode=UseNone,
}}
\usepackage[scale=0.80]{geometry}
\usepackage{import}

% dados pessoais
\name{[First]}{[Last]}
% Convenção brasileira: apenas cidade e estado, nunca rua e número - o mesmo PDF
% circula por muitas empresas. Se não houver o que listar, APAGUE a linha inteira.
% \address{}{}{} falha com "There's no line here to end" em toda versão do moderncv.
\address{[Cidade, UF]}{}{}
\phone[mobile]{[+55 XX XXXXX-XXXX]}
\email{[seu.email@exemplo.com]}
\extrainfo{\href{[https://linkedin.com/in/seu-perfil]}{LinkedIn}, \href{[https://github.com/seu-usuario]}{GitHub}}
% Sem \photo: foto atrapalha a leitura por ATS e não acrescenta nada.
% Sem CPF, RG, estado civil, idade ou data de nascimento - dado sensível sob a LGPD.

\begin{document}

\makecvtitle

% ============================================================
%     RESUMO PROFISSIONAL
% ============================================================

\vspace{6pt}
\small{[Escreva um resumo de 3-5 linhas direcionado à vaga em questão. Foque no que te torna singularmente qualificado. Exemplo: "Cientista de dados com doutorado em [área] e X anos de experiência em indústria. Combina domínio profundo de [domínio] com machine learning aplicado e desenvolvimento em Python. Experiência em pipelines de ML ponta a ponta, da ingestão de dados a dashboards para áreas de negócio."]}

% ============================================================
%     COMPETÊNCIAS
% ============================================================

\section{Competências}
\vspace{1pt}
\begin{itemize}

\item \textbf{[Categoria de competência 1]}: [Competências, frameworks e ferramentas específicas. Seja concreto.]

\item \textbf{[Categoria de competência 2]}: [Competências, frameworks e ferramentas.]

\item \textbf{[Categoria de competência 3]}: [Competências, frameworks e ferramentas.]

\item \textbf{[Categoria de competência 4]}: [Domínio, métodos, abordagens.]

\item \textbf{[Categoria de competência 5]}: [Ferramentas, plataformas, softwares.]

\end{itemize}

% ============================================================
%     EXPERIÊNCIA PROFISSIONAL
% ============================================================

\section{Experiência Profissional}
\vspace{3pt}
\begin{itemize}

% --- Cargo mais recente ---
\item{\cventry{[AAAA-Atual]}{[Cargo]}{[Empresa]}{[Cidade, UF]}{}{\vspace{1pt}
\begin{itemize}
    \item {[Realização ou responsabilidade 1 - seja específico, use números quando possível]}
    \item {[Realização ou responsabilidade 2]}
    \item {[Realização ou responsabilidade 3]}
    \item {[Realização ou responsabilidade 4]}
\end{itemize}}}

\vspace{3pt}

% --- Cargo anterior ---
\item{\cventry{[AAAA-AAAA]}{[Cargo]}{[Empresa]}{[Cidade, UF]}{}{\vspace{1pt}
\begin{itemize}
    \item {[Realização ou responsabilidade 1]}
    \item {[Realização ou responsabilidade 2]}
    \item {[Realização ou responsabilidade 3]}
\end{itemize}}}

\vspace{3pt}

% --- Cargo mais antigo ---
\item{\cventry{[AAAA-AAAA]}{[Cargo]}{[Empresa]}{[Cidade, UF]}{}{\vspace{1pt}
\begin{itemize}
    \item {[Realização ou responsabilidade 1]}
    \item {[Realização ou responsabilidade 2]}
\end{itemize}}}

\end{itemize}

% ============================================================
%     FORMAÇÃO ACADÊMICA
% ============================================================

\section{Formação Acadêmica}
\vspace{1pt}
\begin{itemize}

\item{\cventry{[AAAA-AAAA]}{[Grau] em [Área]}{[Instituição]}{[Cidade, UF]}{}{\vspace{1pt}
TCC/Dissertação: ``[Título do trabalho].'' [Breve descrição do foco da pesquisa.]
}}

\vspace{3pt}

\item{\cventry{[AAAA-AAAA]}{[Grau] em [Área]}{[Instituição]}{[Cidade, UF]}{}{\vspace{1pt}
[Breve descrição ou temas principais.]
}}

\end{itemize}

% ============================================================
%     IDIOMAS
% ============================================================

\section{Idiomas}
\vspace{1pt}
\begin{itemize}
\item {Português (nativo), [Idioma 2] ([nível]), [Idioma 3] ([nível]).}
\end{itemize}

% ============================================================
%     PUBLICAÇÕES (opcional)
% ============================================================

\section{Publicações}
\vspace{1pt}
\begin{itemize}
\item {[Autor(es)] ([Ano]). [Título]. [Periódico/Congresso]. \href{[URL_DOI]}{link DOI}}
\end{itemize}

% ============================================================
%     PRÊMIOS E RECONHECIMENTOS (opcional)
% ============================================================

\section{Prêmios e Reconhecimentos}
\vspace{1pt}
\begin{itemize}
\item{\textbf{[Nome do prêmio]} -- [Evento/Organização] ([Ano]).}
\end{itemize}

% ============================================================
%     REFERÊNCIAS
% ============================================================

\section{Referências}
\vspace{1pt}
\begin{itemize}
\item{Disponíveis mediante solicitação.}
\end{itemize}

\end{document}
